% ============================================================
%  chapter07.tex — 第七章 环面上的常微系统
%  对应原书第七章前四节。
% ============================================================

\setchapterimage{}
\setchapterintro{}
\chapter{环面上的常微系统}

\newcommand{\chSevenFigurePlaceholder}[1]{%
  \begin{center}
    \fbox{\parbox{0.72\textwidth}{\centering\vspace{2.4cm}图~#1\\（图片占位）\vspace{2.4cm}}}
  \end{center}%
}

% ============================================================
%  §1  引言
% ============================================================
\section{引言}

本章讨论环面上的常微系统的经典结果，主要内容参考文献~[1]、[2]。设一阶方程组
\[
  \dfrac{\mathrm d\varphi}{\mathrm dt}=F(\varphi,\theta),\qquad
  \dfrac{\mathrm d\theta}{\mathrm dt}=G(\varphi,\theta), \tag{1.1}
\]
其中函数 $F,G$ 在平面上连续，并且
\[
  F(\varphi+1,\theta)=F(\varphi,\theta)=F(\varphi,\theta+1),
\]
\[
  G(\varphi+1,\theta)=G(\varphi,\theta)=G(\varphi,\theta+1).
\]
因此 $F,G$ 为周期函数，系统~(1.1) 在 $(-\infty,+\infty)$ 上有定义。

若把 $(\varphi,\theta)$ 平面内单位正方形的对边粘合，那么便产生一个环面 $T^2$；
系统~(1.1) 可以解释为环面 $T^2$ 上的一个常微系统。为便于直观起见，在三维 Euclid
空间内可把 $T^2$ 写为
\[
  T^2:\quad \left(\sqrt{x^2+y^2}-a\right)^2+z^2=b^2,
\]
其中 $a>b>0$。它的点可由参数表示为
\[
  x=(a+b\cos2\pi\theta)\cos2\pi\varphi,\qquad
  y=(a+b\cos2\pi\theta)\sin2\pi\varphi,
\]
\[
  z=b\sin2\pi\theta,\qquad
  0\leq\varphi<1,\quad 0\leq\theta<1,\quad 0<b<a. \tag{1.1a}
\]

\chSevenFigurePlaceholder{7.1}

图~7.1 所表示的曲线在平面上看似不闭合，但在环面上两条对边相同，故对应于
$T^2$ 上的一条连续曲线。我们把微分系统~(1.1) 规定在环面上，再在环面 $T^2$ 上
作相平面的话，可以研究它的周期解和轨线的长期性态。

下面常作一个重要假定：
\[
  F(\varphi,\theta)>0. \tag{1.2}
\]
此时 $\varphi$ 沿每一条轨线严格增加，可取 $\varphi$ 为自变量，把~(1.1) 化为
\[
  \dfrac{\mathrm d\theta}{\mathrm d\varphi}
  =A(\varphi,\theta):=\dfrac{G(\varphi,\theta)}{F(\varphi,\theta)}. \tag{1.3}
\]
若 $\theta=\theta(\varphi;\varphi_0,\theta_0)$ 是~(1.3) 的解，则它在 $T^2$ 上的像
就是系统~(1.1) 的一条轨线，方向为 $\varphi$ 增加的方向。

\begin{example}{}{ch7-intro-example}
设
\[
  \dfrac{\mathrm d\varphi}{\mathrm dt}=1,\qquad
  \dfrac{\mathrm d\theta}{\mathrm dt}=\mu,
\]
其中 $\mu$ 为常数，则
\[
  \theta=\mu\varphi+\theta_0.
\]
若 $\mu=m/n$，其中 $m,n$ 互素且 $n>0$，则每一条轨线均为闭轨线；若 $\mu$ 为无理数，
则每一条轨线在环面上稠密。
\end{example}

% ============================================================
%  §2  旋转数
% ============================================================
\section{旋转数}

设常微分方程
\[
  \dfrac{\mathrm d\theta}{\mathrm d\varphi}=A(\varphi,\theta) \tag{2.1}
\]
满足下列条件：
\[
  A(\varphi,\theta)\ \text{连续且有界},\qquad
  A(\varphi+1,\theta)=A(\varphi,\theta)=A(\varphi,\theta+1),
\]
并设初值问题的解唯一。令
\[
  \theta=\psi(\varphi;\varphi_0,\theta_0)
\]
表示~(2.1) 满足初始条件 $\theta_0$ 的解。由于方程关于两个变量均以 $1$ 为周期，
有
\[
  \psi(\varphi+1;\varphi_0+1,\theta_0)
  =\psi(\varphi;\varphi_0,\theta_0),
\]
\[
  \psi(\varphi;\varphi_0,\theta_0+1)
  =\psi(\varphi;\varphi_0,\theta_0)+1. \tag{2.3}
\]

取截面 $\varphi=0$，并令
\[
  \Phi(\theta_0)=\psi(1;0,\theta_0). \tag{2.2}
\]
于是 $\Phi$ 是实直线上的严格单调连续映射，且
\[
  \Phi(\theta+1)=\Phi(\theta)+1.
\]
记 $\Phi^n$ 为它的 $n$ 次迭代，则 $\Phi^n(\theta_0)$ 对应于解在 $\varphi=n$ 时的
$\theta$ 坐标。

\begin{theorem}{}{ch7-rotation-number}
对任意满足上述条件的方程~(2.1)，极限
\[
  \mu=\lim_{|n|\to\infty}\dfrac{\Phi^n(\theta_0)}{n} \tag{2.4}
\]
存在，且与初值 $\theta_0$ 无关。数 $\mu$ 称为方程~(2.1) 或映射 $\Phi$ 的
\textbf{旋转数}。
\end{theorem}

证明的要点如下。由 $\Phi$ 的单调性，若 $m<\theta_0<m+1$，便有
\[
  \Phi^n(m)<\Phi^n(\theta_0)<\Phi^n(m+1).
\]
又由周期关系可将任意两个初值的迭代差控制在一个固定常数内。令
$\theta_n=\Phi^n(\theta_0)$，则序列 $\theta_n-\theta_0$ 具有近似可加性；由上下逼近可得
极限~(2.4) 存在。

旋转数的有理性同周期点的存在密切相关。若存在整数 $m,n$，$n>0$，使
\[
  \Phi^n(\theta_0)=\theta_0+m,
\]
则相应轨线为闭轨线，且 $\mu=m/n$。反之，若 $\mu$ 是有理数，适当的迭代映射
$\Phi^n-m$ 有不动点；故在本节假定下，方程有周期轨线。

% ============================================================
%  §3  极限点集
% ============================================================
\section{极限点集}

类似于平面动力系统的极限点集，设 $S$ 是环面 $T^2$ 上的一条轨线，记
\[
  S^+=\{P(t):0\leq t<+\infty\},\qquad
  S^-=\{P(t):-\infty<t\leq0\}.
\]
若 $P\in T^2$，$S_P^+$ 表示从点 $P$ 出发的正半轨线。定义
\[
  S'_P=\bigcap_{t\geq0}\overline{S_P(t,+\infty)}
\]
为 $P$ 的正极限点集；负极限点集定义相同。

\begin{lemma}{}{ch7-limit-invariant}
对任意 $P\in T^2$，集合 $S'_P$ 非空、闭且不变；若 $Q\in S'_P$，则从 $Q$ 出发的
全轨线完全包含于 $S'_P$。
\end{lemma}

事实上，$T^2$ 为紧致空间，故正半轨线任取一点列总有收敛子列。若
$P_i=P(t_i)\to Q$，$t_i\to+\infty$，则由解对初值的连续依赖性，对任意固定 $s$ 有
\[
  P(t_i+s)\longrightarrow Q(s),
\]
这就说明 $Q(s)\in S'_P$。

为了描述可能出现的极限点集，设 $P_0,P_1,\ldots,P_m$ 是环面上一条横截线上的点，
并且这些点在截面上按同一方向排列。设 $C$ 是由相邻轨线段所围成的带状区域。

\begin{lemma}{}{ch7-limit-order}
若 $P_i$ 的正半轨线两两不相交，且在横截线上的顺序保持，则这些轨线的极限点集相同
或彼此不交。特别地，若其中一条轨线的正极限点集为整个环面，则所有相邻轨线也有
同样的正极限点集。
\end{lemma}

\begin{theorem}{}{ch7-limit-classification}
设 $S$ 是环面 $T^2$ 上方程~(2.1) 的一条轨线，且系统在 $T^2$ 上无奇点。则 $S$ 的
正极限点集只有下列可能：\newline
\noindent (i) 一条闭轨线；\newline
\noindent (ii) 整个环面 $T^2$；\newline
\noindent (iii) 一个由闭轨线围成的不变带状区域的边界。
\end{theorem}

这里，若 $S'_P$ 不含周期轨线，则由截面映射的单调性和紧致性可知，$S'_P$ 中任一点
的轨线均在其中稠密；在通常的非奇异情形，便得到 $S'_P=T^2$。

% ============================================================
%  §4  各态经历
% ============================================================
\section{各态经历}

当系统~(2.1) 是解析的情形，Poincar\'e 曾经猜想不会出现奇异的情形；A.~Denjoy 在
1932 年证明了 Poincar\'e 的猜想。为叙述该结果，设
\[
  A(\varphi,\theta)\ \text{连续，且}\quad A(\varphi,\theta)>0,
\]
并且 $A$ 关于 $(\varphi,\theta)$ 具有足够的光滑性。令截面映射仍为
\[
  \Phi(\theta_0)=\psi(1;0,\theta_0).
\]

\begin{theorem}{Denjoy}{ch7-denjoy}
设 $\Phi$ 是圆周上的保向同胚，且 $\Phi'$ 绝对连续，$\Phi''$ 可积。若 $\Phi$ 的
旋转数为无理数，则 $\Phi$ 的每一条轨道在圆周上稠密。等价地，系统~(2.1) 的每一条
轨线在环面 $T^2$ 上稠密。
\end{theorem}

定理说明，在上述光滑性条件下，环面系统具有各态经历性。为了证明它，取一点
$P_0$ 及其迭代点
\[
  P_n=\Phi^n(P_0),\qquad n=0,\pm1,\pm2,\ldots . \tag{4.1}
\]
当旋转数无理时，这些点两两不同。若它们不稠密，则其闭包的补集中存在一个开区间；
Denjoy 的论证通过考察该区间在 $\Phi$ 的所有迭代下的像，得到矛盾。

令
\[
  \omega(\theta)=\dfrac{\partial\psi}{\partial\theta}(1;0,\theta)>0.
\]
在适当的正则性条件下，$\log\omega(\theta)$ 有有界变差，并且对相邻迭代点满足
\[
  \left|\log\dfrac{\mathrm d\Phi^n(\theta)}{\mathrm d\theta}\right|\leq V, \tag{4.2}
\]
其中 $V$ 为与 $n$ 无关的常数。这个有界畸变估计保证：若某一开区间的全部迭代像
两两不交，则其长度总和不能同时满足旋转数为无理数所要求的逼近性质，故矛盾。

\begin{theorem}{}{ch7-ergodic}
设 $A(\varphi,\theta)$ 连续，并且其一阶偏导数连续。又设 $A(\varphi,\theta)>0$，且
对应的旋转数为无理数。则方程~(2.1) 的每一条轨线在环面 $T^2$ 上稠密。
\end{theorem}

若 $\Phi$ 的光滑性不足，上述结论可以失效；第 5 节将给出奇异情况的例子，第 6 节
介绍 Schweitzer 的反例。

% ============================================================
%  §5  奇异情况举例
% ============================================================
\section{奇异情况举例}

在这一节中，我们举出奇异情形的例子。设圆周 $C$ 上给定一个 Cantor 集合 $F$，
设它的补集是可数个两两不交的开区间
\[
  (a_\beta,b_\beta),\qquad \beta=1,2,\ldots .
\]
设 $\mu_k$ 为正整数，并把单位圆周上的点按所选定的序列编号。可选取一列点
\[
  0,\ \mu_1,\ -\mu_2,\ 2\mu_1-2\mu_2,\ldots,\ k\mu_1-k\mu_2,\ldots \tag{5.1}
\]
使所有这些点均在 $F$ 的同一稠密子集中。由此把各个区间 $(a_\beta,b_\beta)$ 之间的
对应关系逐步规定出来。

更具体地，在每个补区间上作一个严格递增的同胚，把右端点映到相应区间的右端点，
左端点映到相应区间的左端点；在 Cantor 集合 $F$ 上把映射规定为平移。这样得到圆周
上的保向同胚 $T_1$，它满足
\[
  T_1(\xi+1)=T_1(\xi)+1.
\]
其轨道可在 $F$ 的补区间之间游走，却不必在圆周上稠密。适当选取各区间上的函数，
可使 $T_1$ 连续可微，且
\[
  0<\varepsilon<T_1'(\xi)<T_1'(\eta)<T_1(\xi)+1
\]
在相应的区间上成立。

把 $T_1$ 看作截面 $\varphi=0$ 上的 Poincar\'e 映射。于是可在带状区域
\[
  0\leq\varphi\leq1
\]
上构造一个非奇异向量场 $(P,Q)$，使得
\[
  \frac{\mathrm d\varphi}{\mathrm dt}=1,\qquad
  \frac{\mathrm d\theta}{\mathrm dt}=T_1(\theta)-\theta. \tag{5.2}
\]
在两条边界上按
\[
  (\varphi,\theta)\sim(\varphi+1,T_1(\theta))
\]
粘合，便得到环面上的常微系统。该系统的旋转数为无理数，但存在一个非空的闭不变集
$F$，且从补区间出发的轨线永远不能穿入 $F$ 的内部。因此，若截面映射的光滑性不足，
第 4 节的各态经历结论可以失效。

% ============================================================
%  §6  介绍 Schweitzer 之例
% ============================================================
\section{介绍Schweitzer之例}

在 1950 年，Seifert 提出下面的问题：设 $\mathscr S^3$ 是三维 Euclid 空间内的一个
光滑闭流形，在 $\mathscr S^3$ 上给定一个非奇异的 $C^r$ 向量场；它是否必有一条
闭轨线？这个问题是 Poincar\'e--Bendixson 定理在高维空间中的一个推广问题。1974 年
Schweitzer 发表了一个否定的例子。

设 $N$ 是一个有边界的二维流形，在 $N\times D^1$ 上取坐标 $(X,Z)$，并设
$X_0=(0,0)$。令
\[
  B=N\times D^1,
\]
其中 $D^1$ 为闭区间。先在 $N$ 上构造一个具有一个孤立奇点 $X_0$ 的平面向量场 $X$，
使得 $X$ 在 $N\setminus\{X_0\}$ 上非零，并满足：
\[
  \text{从 }(n-1)\text{ 到 }n\text{ 的每一条轨线均穿过指定的横截面，}
\]
\[
  \text{且在 }X_0\text{ 附近的轨线不闭合。} \tag{6.1}
\]
然后把 $X$ 延拓到 $N\times D^1$，并在靠近边界的区域将其同一个竖直向量场拼接。

令
\[
  X_1=N\times I-(1-\varphi)X_0+\varphi Z,
\]
其中 $\varphi$ 是适当的光滑截断函数。通过选择 $\varphi$，可使 $X_1$ 在边界邻域内
同标准竖直向量场相同，而在内部保留原有的非周期性态。

下面说明 Schweitzer 例的拓扑构造。在 $\mathscr S^3$ 内可放置一个非奇异的
$C^r$ 向量场 $Y$，使得在其附近存在一个实心环面 $E$。取
\[
  B=N\times D^1,
\]
并把 $B$ 沿两个横截面嵌入 $E$ 中。经过一个同胚变形，原来的一个实心环面被替换为
两个相互套合的实心环面；在边界上向量场仍保持横截。

\chSevenFigurePlaceholder{7.2}
\chSevenFigurePlaceholder{7.3}

在所得流形中，任何进入 $B$ 的轨线都按构造在有限时间后离开 $B$，但它不可能回到
原来的截面而形成闭轨线；而不进入 $B$ 的轨线则由原向量场的性质也不能形成闭轨线。
于是经粘合得到的 $\mathscr S^3$ 上的非奇异 $C^r$ 向量场没有闭轨线。

\chSevenFigurePlaceholder{7.4}

图~7.4 表示在圆柱形邻域内的局部坐标。取中心轴上的点列
\[
  P_0,P_1,P_2,\ldots
\]
作为横截面上的基点；向量场的流把相邻截面之间的轨线按预先给定的方式连接。由此，
所有可能的回归轨线均被引入插入区，因而不能闭合。这就是 Schweitzer 反例的基本
思想。

% ============================================================
%  习题和参考文献
% ============================================================
\section*{习题和参考文献}
\subsection*{习题}

\begin{enumerate}
  \item 设映射 $\Phi$ 有周期。证明方程~(2.1) 的一条轨线满足相应的周期条件，当且仅当
  它在截面映射 $\Phi$ 下为周期点。
  \item 下面所有函数均为 $C^1$ 连续。设周期为 $1$，并且能证明以 $0$ 为初值的
  \[
    \frac{\mathrm d\theta}{\mathrm d\varphi}
    =a\varphi+\sin2\pi\theta
  \]
  的旋转数存在。再讨论
  \[
    \frac{\mathrm d\theta}{\mathrm d\varphi}
    =\varphi\sin2\pi\theta+\varepsilon g(\theta,\varphi)
  \]
  的旋转数。
  \item 设 $\omega=\omega(\varphi,\theta)>0$，证明存在连续函数
  $h(\varphi,\theta)$，使相应截面映射具有指定的无理旋转数；进一步在
  $h$ 满足 Lipschitz 条件时，利用第 4 节的结果讨论环面 $T^2$ 上各态经历轨线的存在。
\end{enumerate}

\subsection*{本章参考文献}
\begin{enumerate}
  \item Coddington, E.~A. and Levinson, N., \textit{Theory of Ordinary Differential Equations},
  McGraw--Hill, 1955.
  \item Hale, J.~K., \textit{Ordinary Differential Equations}, Wiley--Interscience, 1969;
  2nd edition, Krieger Publ. Co., 1980.
  \item Schweitzer, P.~A., Counterexample to the Seifert conjecture and opening the closed
  orbits of flows, \textit{Ann. of Math.}, 100 (1974), 386--388.
  \item 史耐德--孔德，微分方程常微分解构造法，\textit{数学学报}，1977，136--148。
\end{enumerate}

\endinput
